\section{Meshes, Boundary Tags, and Case Inputs}
\label{sec:mesh}

\subsection{CGNS input}

Meshes are read through the CGNS mid-level library. The reader walks every base
and every zone, accepts \texttt{Unstructured} zones only, and reads the
coordinate arrays \texttt{CoordinateX}/\texttt{CoordinateY}. The two supplied meshes
differ here, which is why the reader cannot assume either layout:
\texttt{NACA0012\_H2.cgns} declares \texttt{PhysicalDimension} $=3$ and carries a
nominally zero \texttt{CoordinateZ}, which is read and discarded, while
\texttt{CylinderB1.cgns} declares \texttt{PhysicalDimension} $=2$ and has no $z$
coordinate at all. Both declare \texttt{CellDimension} $=2$. The solver logs the
declared dimensions for each mesh it opens. Element sections are scanned
for the supported two-dimensional element types \texttt{TRI\_3} and
\texttt{QUAD\_4}, which become cells, and \texttt{BAR\_2}, which becomes a
boundary edge. No structured-grid or single-element-type assumption is made
anywhere: the NACA mesh is genuinely mixed, with $\nacaTri$ triangles and
$\nacaQuad$ quadrilaterals.

Two CGNS subtleties required specific handling and are worth recording:

\begin{itemize}
  \item \textbf{Boundary conditions.} The \texttt{BC\_t} nodes store a
  \texttt{PointRange} with \texttt{GridLocation=EdgeCenter}. These indices are
  \emph{element} indices, not node indices, and they coincide with the
  \texttt{ElementRange} of the identically named \texttt{BAR\_2} section. The
  reader therefore maps each boundary condition to its edge section by range
  containment, then tags the resulting mesh edges with the family name. Family
  names are carried as strings and matched against the
  \texttt{boundary\_conditions} dictionary of the case file, so
  \texttt{bc-2}/\texttt{bc-4} and \texttt{WALL}/\texttt{FAR} are handled by the
  same code path. An unmapped or unknown family is a fatal, clearly reported
  error.
  \item \textbf{Multi-zone connectivity.} The cylinder mesh has two zones. Its
  interfaces are stored as \texttt{GridConnectivity\_t} of type
  \texttt{Abutting1to1}, \emph{not} as \texttt{GridConnectivity1to1\_t}: a
  \texttt{cg\_n1to1} query returns zero. The reader therefore uses
  \texttt{cg\_nconns}/\texttt{cg\_conn\_info}/\texttt{cg\_conn\_read} and
  interprets the \texttt{PointList}/\texttt{PointListDonor} pairs as one-based
  zone-local \emph{vertex} indices (the SIDS default location). Paired nodes are
  merged into single global nodes; the supplied interface point lists are
  bit-identical, so the maximum merge distance is exactly zero. After merging,
  the cylinder mesh has $\cylNodes$ nodes and $\cylCells$ cells.
\end{itemize}

\subsection{Face extraction and geometry}

Cells are converted to an internal polygon representation and a face list is
built by hashing each edge's sorted node pair. An edge seen by two cells becomes
an interior face with a left and a right cell; an edge seen once becomes a
boundary face carrying its family tag. This yields the cell--face adjacency
graph used for reconstruction, for the gradient stencils and, in its dual form,
for partitioning.

Cell area and centroid come from the shoelace formula over the vertex loop,
\begin{equation}
  A=\frac12\sum_i \left(x_i y_{i+1}-x_{i+1}y_i\right),\qquad
  \mathbf{x}_c=\frac{1}{6A}\sum_i (\mathbf{x}_i+\mathbf{x}_{i+1})
  \left(x_i y_{i+1}-x_{i+1}y_i\right),
  \label{eq:shoelace}
\end{equation}
which is exact for both triangles and quadrilaterals and needs no subdivision.
Cell node ordering is \emph{normalised} rather than assumed: any cell whose signed
area comes out negative has its node list reversed, and the number of reorientations
is logged. Only a zero or degenerate area is fatal, since that cannot be repaired by
reordering. For both supplied meshes the reorientation count is zero --- every cell
is already stored counter-clockwise --- so the normalisation is a no-op here, but the
reader does not depend on that being true. In two dimensions a
face is a straight segment from $\mathbf{x}_a$ to $\mathbf{x}_b$, so its length
$\Delta s$ (the two-dimensional analogue of face area) and unit normal are
\begin{equation}
  \Delta s=\|\mathbf{x}_b-\mathbf{x}_a\|,\qquad
  \mathbf{n}=\frac{1}{\Delta s}\left(y_b-y_a,\ -(x_b-x_a)\right),
\end{equation}
with the face centroid at the segment midpoint. Each face normal is then
oriented outward from its left cell by testing the sign of
$\mathbf{n}\cdot(\mathbf{x}_f-\mathbf{x}_{c,\mathrm{left}})$ and flipping when
required. On a boundary face the left cell is by construction the adjacent fluid
cell, so \textbf{the boundary normal points out of the fluid and into the body}.
That convention matters for the force integration in \cref{sec:forces} and is
the single easiest sign error to make in a solver of this kind.

\subsection{Mesh and case summary}

\Cref{tab:meshes} summarises both meshes. The wall-boundary counts are what the
surface output and force integration act on: $\nacaWallEdges$ edges around the
airfoil and $\cylWallEdges$ around the cylinder.

\begin{table}[htbp]
  \centering
  \caption{Supplied meshes as read and preprocessed by the solver. The cylinder
  figures are after the two zones have been merged across their
  \texttt{Abutting1to1} interface.}
  \label{tab:meshes}
  \small
  \begin{tabular}{lrr}
    \toprule
    Quantity & NACA0012\_H2 & CylinderB1 \\
    \midrule
    CGNS zones & 1 & 2 (merged) \\
    Nodes & $\nacaNodes$ & $\cylNodes$ \\
    Cells & $\nacaCells$ & $\cylCells$ \\
    \quad triangles & $\nacaTri$ & mixed \\
    \quad quadrilaterals & $\nacaQuad$ & mixed \\
    Faces & $\nacaFaces$ & $\cylFaces$ \\
    Boundary faces & $\nacaBoundaryFaces$ & $\cylBoundaryFaces$ \\
    Wall-family edges & $\nacaWallEdges$ (\texttt{bc-4}) & $\cylWallEdges$ (\texttt{WALL}) \\
    Farfield edges & $\nacaFarEdges$ (\texttt{bc-2}) & $\cylFarEdges$ (\texttt{FAR}) \\
    Farfield extent & $\approx 80\,c$ & $200\,D$ \\
    Body extent & $x\in[0,1.005]$, $t=0.120$ & $r=0.5$ exactly \\
    \bottomrule
  \end{tabular}
\end{table}

A note on the airfoil geometry: the \texttt{bc-4} family spans $x\in[0,1.005]$
with a maximum thickness of $0.120$, consistent with a NACA0012 section of unit
chord. Its identification as the body, rather than \texttt{bc-2} at radius
$\approx 80$, comes from the case-file mapping and not from geometry.
